\section{Introduction}
\label{sec:introduction}

An inertial navigation system (INS) integrates specific force and angular rate into position, and its error grows without bound unless an external position or velocity reference arrests it. On most platforms that reference is GNSS, whose low received power makes it easy to deny by jamming, to corrupt by spoofing, and to lose in urban canyons, under foliage, indoors, and underwater \citep{humphreys2008, divis2013gps}. How quickly the INS error grows once GNSS is lost is set by the grade of the inertial sensor. A navigation-grade system coasts for tens of minutes before accumulating 100~m of error; the consumer-grade micro-electro-mechanical systems (MEMS) IMU found in a smartphone, a small drone, or a low-cost autonomous vehicle crosses the same threshold in under a minute \citep{groves_ch14}. For the platforms that carry such sensors, which vastly outnumber those carrying navigation-grade instruments, an alternative aiding source must be available continuously, everywhere the platform operates, and at a cost commensurate with the platform itself.

Map-matching against the Earth's geophysical fields is one of the few alternatives that is present everywhere and at all times without depending on infrastructure. Terrain- and bathymetry-aided navigation are the most mature members of that family \citep{baker1977terrain, williams2001towards_tan, donovan2012}, but they require an active emission. Gravity and magnetic anomalies are sensed passively, and anomaly-aided navigation has been demonstrated with gravimeters on submarines and ground vehicles \citep{jircitano1991gravity, kamgar1999vehicle} and with survey magnetometers on aircraft \citep{canciani2016absolute, canciani2017airborne}. Every such demonstration, however, pairs a tactical- or navigation-grade INS with a specialized survey instrument. Whether the technique survives the replacement of both halves of that pairing, a MEMS IMU in place of the high-grade inertial system and commodity sensing in place of the survey instrument, has not been established.

This paper addresses that question with three estimators on one body of data. An extended Kalman filter (EKF), an unscented Kalman filter (UKF), and a Rao-Blackwellized particle filter (RBPF) in the form of \citet{canciani2017airborne} are each aided with gravity anomaly, magnetic anomaly, or both, and each aided run is paired with the same filter's unaided run on the same drive. All three filters share one strapdown mechanization and one seeded measurement stream, so that the anomaly measurement is the only difference within a pair. The drives are the 27 trajectories of the MEMS-Nav dataset, \numdatasetkm~km and \numdatasethours~h of smartphone recordings on public roads in the Mid-Atlantic United States \citep{mems-nav-dataset}. GNSS is degraded to one fix per minute with correlated position and velocity errors, a regime in which the unaided EKF's median horizontal error is \numekfdegradedmedian~m against \numekffullmedian~m with full GNSS. The anomaly observations come from two sources: the gravity and magnetic readings that the smartphone itself records, and, as a control, simulated readings of two low-cost dedicated sensors substituted for them with everything else held fixed.

\subsection{Relation to Earlier Work by the Authors}

Preliminary versions of this work were presented at the 2026 ION International Technical Meeting, with a UKF \citep{anom-ukf}, and at the 2026 ION Pacific PNT Meeting, with an RBPF \citep{anom-pf}. Both reported that the smartphone's gravity and magnetic readings reduced navigation error, by tens to hundreds of meters for the UKF and by more than 1{,}500~km for the RBPF. Those results do not hold. An audit of the software found that the anomaly observation was formed with unit errors that left it unrelated to the maps, and that the comparisons were made against baselines that had themselves diverged. \Cref{sec:corrections} sets out each defect and its effect. This paper replaces the conference results in full: it uses the corrected measurement models, a larger dataset, a common GNSS degradation for every run, a paired statistical design, and a third estimator. Its conclusions about the smartphone sensors are the opposite of those reported earlier.

\subsection{Contributions}

The contributions of this paper are as follows.
\begin{enumerate}
  \item A paired evaluation of EKF, UKF, and RBPF anomaly aiding of a consumer-grade MEMS INS on 27 ground-vehicle drives, in which every filter receives an identical, reproducible measurement stream and each aided run is scored against its own unaided counterpart.
  \item A corrected formulation of the gravity and magnetic anomaly observations, with an error model measured from the data against the maps, and an explicit account of the defects that invalidated the authors' earlier conference results.
  \item A null result for the smartphone's own sensors and its mechanism: the operating system pins the reported gravity norm to a device constant, and the vehicle's own magnetic field dominates the magnetometer, so both observations have signal-to-noise ratios far below one against the maps.
  \item A controlled test with simulated dedicated sensors showing that a \$25 magnetometer makes magnetic aiding consistently beneficial for both Kalman filters with the same consumer-grade IMU, whereas a \$77 MEMS accelerometer used as a gravimeter does not, and that the size of the benefit is consistent with a per-update localization bound.
  \item A negative result for the particle filter, which loses track under once-per-minute GNSS with or without aiding, with a diagnosis of the failure.
  \item An open-source implementation \citep{strapdown-rs} and dataset with which every number in this paper can be regenerated.
\end{enumerate}

The remainder of the paper is organized as follows. \Cref{sec:corrections} documents the corrections to the conference results. \Cref{sec:navigation} describes the shared mechanization, the GNSS degradation model, and the three filters. \Cref{sec:geophysical} develops the anomaly measurement models and their error models. \Cref{sec:experiment} describes the data and the experimental design, and \Cref{sec:results} the results. \Cref{sec:discussion} discusses the sensor requirement that the results imply, the behavior of the particle filter, and the limitations of the study, and \Cref{sec:conclusion} concludes.
